% GDZ PPN236018647,native575;HKUST576.
\clearpage\SourceLeaf{575}\ArithmeticHead{546}{L}{TAGEBUCH MIT ERLÄUTERUNGEN.}
\NumberedItem{105}
\indent Theoriam quantitatum transcendentium:
\[
\int\frac{dx}{\surd(1-\alpha xx)(1-\beta xx)}
\]
\noindent ad summam universalitatem perduximus.
\par\noindent\hfill Brunov[ici, 1800] Mai. 6.

\medskip\begingroup\small
\makeatletter\tagsleft@true\let\veqno\@@leqno\makeatother
\indent Die Aufdeckung des Zusammenhangs zwischen dem agM. und dem vollständigen elliptischen Integral erster Gattung, sowie die in dem besondern Falle der lemniskatischen Funktionen gewonnenen Einsichten, führten \textsc{Gauss} zu der Erkenntnis, daß die Größen
\begin{equation*}
\tag*{(*)}\frac c{M(1,c)},\quad\frac c{M(1,s)}\qquad(s=\sqrt{1-c^2})
\end{equation*}
\noindent (siehe den art. [7.], oben S. 190, aus Scheda Ac, S. 13) für das elliptische Integral mit dem Modul \(s\) die analoge Rolle spielen, wie die Größe
\[
\frac1{M(\sqrt2,1)}=\frac\varpi\pi\mathpunct{\raisebox{.35ex}{.}}
\]
\noindent auf die sie sich für \(s=\sqrt{\frac12}\) reduzieren, für die Lemniskate. Daß \textsc{Gauss}’ Auffassung dieser Größen (der Periodizitätsmoduln) dem neuern Standpunkte sehr nahe gewesen sein muß, zeigt die folgende Aufzeichnung, die sich auf einem Zettel ohne Datum findet:
\par\begingroup\normalsize
\indent Der Radicalfehler, woran meine bisherigen Bestrebungen, den Geist der elliptischen Function zu verkörpern, gescheitert sind, scheint \textls[120]{der} zu sein, dass ich dem Integral
\begin{equation*}
\tag*{(**)}\int\frac{d\varphi}{\surd(1-e^2\sin\varphi^2)}
\end{equation*}
\noindent die Bedeutung als Ausdruck eines endlichen Theils der Kugelfläche habe unterlegen wollen, während es wahrscheinlich nur einen unendlich schmalen Kugelsector ausdrücken soll.
\par\endgroup
\indent Offenbar bedeutet hier die »Kugelfläche« den Ort der komplexen Veränderlichen, der »endliche Theil«, dessen »Ausdruck« das Integral (**) sein sollte, den Fundamentalbereich oder das Periodenparallelogramm; wenn dann \textsc{Gauss} von »einem unendlich schmalen Kugelsektor« spricht, so zeigt dies, daß er damals, als er diesen Zettel schrieb, es für wahrscheinlich hielt, daß man die reellen Größen (*) selbst, oder genauer (vergl. oben S. 194 Gln. [1])
\[
\frac{2\pi}{M(1,\sqrt{1-e^2})},\quad\frac{2\pi}{M(1,e)}
\]
\noindent als die Periodizitätsmoduln des Integrals (**), für \(e\) reell und kleiner als Eins, anzusehen habe. Am 6. Mai hatte er sich zur völligen Klarheit über diese Verhältnisse durchgerungen, und damit war der Weg frei für weitere Fortschritte. Vergl. den Abschnitt IV des Aufsatzes »Über \textsc{Gauss}’ Arbeiten zur Funktionentheorie«, Werke X\,2.
\par\noindent\hfill\textsc{Klein.\quad Schlesinger.}
\par\endgroup
